\documentclass[10pt]{article}
\usepackage[T1]{fontenc}
\usepackage[utf8]{inputenc}
\usepackage[french]{babel}
\usepackage{amsmath,amssymb,mathrsfs,mathpazo}
\usepackage{geometry}
\geometry{a4paper,margin=2cm}
\usepackage{graphicx,array,tabularx,tabularray,colortbl,multirow,multicol,booktabs,enumitem,pifont,ulem,fancybox,diagbox,cancel}
\usepackage{pstricks,pstricks-add,pst-plot,pst-tree,pst-node,pst-text,pst-func,pst-math,pst-bspline,pst-3dplot,pst-eucl,pst-solides3d}
\let\pscalculate\undefined
\usepackage{pst-bezier}
\usepackage{tikz,pgfplots,tkz-tab}
\usetikzlibrary{arrows,arrows.meta,patterns,calc,positioning,trees,decorations.markings,shapes.geometric,shapes.misc,shapes.arrows,intersections,decorations.pathmorphing,decorations.pathreplacing,backgrounds,shadows,babel}
\pgfplotsset{compat=1.18}
\usepgfplotslibrary{fillbetween}
\usepackage[np]{numprint}
\usepackage{listings,makecell,pas-tableur,twoopt}
\let\barre\undefined
\usepackage[pstricks]{tablvar}
\definecolor{darkgreen}{RGB}{0,100,0}
\providecommand{\red}{\color{red}}
\providecommand{\blue}{\color{blue}}
\providecommand{\green}{\color{green}}
\providecommand{\black}{\color{black}}
\providecommand{\white}{\color{white}}
\providecommand{\gray}{\color{gray}}
\providecommand{\cyan}{\color{cyan}}
\providecommand{\magenta}{\color{magenta}}
\providecommand{\yellow}{\color{yellow}}
\providecommand{\violet}{\color{violet}}
\providecommand{\orange}{\color{orange}}
\providecommand{\euro}{\texteuro}
\providecommand{\EUR}{\texteuro}
\providecommand{\ieme}{\textsuperscript{e}}
\providecommand{\degres}{\ensuremath{^\circ}}
\providecommand{\degree}{\ensuremath{^\circ}}
\providecommand{\vv}[1]{\overrightarrow{#1}}
\providecommand{\Coord}[1]{#1}
\providecommand{\decosix}{}
\providecommand{\starredbullet}{\ensuremath{\bullet}}
\pagestyle{empty}
\setlength{\parindent}{0pt}
\usepackage{adjustbox}\DeclareUnicodeCharacter{2236}{:}\DeclareUnicodeCharacter{2190}{\ensuremath{\leftarrow}}
\providecommand{\R}{\mathbb{R}}
\providecommand{\N}{\mathbb{N}}
\providecommand{\Z}{\mathbb{Z}}
\providecommand{\Q}{\mathbb{Q}}
\providecommand{\D}{\mathbb{D}}
\providecommand{\C}{\mathbb{C}}
\providecommand{\e}{{\mathrm{e}}}
\providecommand{\vect}[1]{\overrightarrow{#1}}
\providecommand{\vectt}[1]{\overrightarrow{\mathrm{#1}}}
\providecommand{\vvt}[1]{\overrightarrow{\mathrm{#1}}}
\providecommand{\np}[2][]{{#2~#1}}
\providecommand{\vv}[1]{{\overrightarrow{#1}}}
\renewcommand{\d}{{\mathrm{d}}}
\providecommand{\degres}{{^\circ}}
\providecommand{\degre}{{^\circ}}
\providecommand{\blue}{\color{blue}}
\providecommand{\red}{\color{red}}
\providecommand{\enskip}{\quad}
\providecommand{\llbracket}{\lbrack\!\lbrack}
\providecommand{\rrbracket}{\rbrack\!\rbrack}
\providecommand{\euro}{{\text{€}}}
\providecommand{\no}{{\text{n°}}}
\providecommand{\interval}[2]{{[#1;#2]}}
\providecommand{\intervFO}[2]{{[#1;#2[}}
\providecommand{\dotfill}{\ldots}

\providecommand{\ds}{\displaystyle}
\providecommand{\textding}[1]{#1}
\providecommand{\Oij}{\(\mathrm{O};\vec{\imath},\vec{\jmath}\)}
\providecommand{\Oijk}{\(\mathrm{O};\vec{\imath},\vec{\jmath},\vec{k}\)}
\providecommand{\Ouv}{\(\mathrm{O};\vec{u},\vec{v}\)}
\providecommand{\intervFF}[2]{\ensuremath{[#1;#2]}}
\renewcommand{\euro}{\ifmmode\text{\texteuro}\else\texteuro\fi}
\DeclareUnicodeCharacter{00B0}{\ifmmode^{\circ}\else\textdegree\fi}
\DeclareUnicodeCharacter{20AC}{\ifmmode\text{\texteuro}\else\texteuro\fi}
% Énoncés : Éducation nationale. Transcription et corrigés : APMEP, Denis Vergès et contributeurs. Découpage et classement : CoopMaths, Sébastien Lozano, Éric Elter et contributeurs. Adaptation de la présentation : Maths Explorer.
% CC BY-SA : https://coopmaths.fr/www/about/
% https://coopmaths.fr/alea/static/e3c/2020/tex/e3c_2020_00_sujet12_3_cor.tex
\begin{document}
\noindent\textit{APMEP / CoopMaths. CC BY-SA. Présentation adaptée par Maths Explorer.}\par\medskip
Dans la figure ci-dessous, on a tracé $C_f$ la courbe représentative d'une fonction $f$ définie et dérivable sur $\mathbb{R}$ ainsi que les tangentes à $C_f$ aux points d'abscisses respectives $-2$, $-1$ et 0.
	
\subsection*{1. Compléter le tableau de valeurs ci-dessous. }
	
	\[
	\begin{array}{|c|c|c|c|}
		\hline	x & -2 & -1 & 0 \\
		\hline
		f(x) & 1 & 1 & 1 \\
		\hline
		f'(x) & 2 & -1 & 2 \\
			\hline
	\end{array}
	\]
	

	\subsection*{2. a. Calculer $f'(x)$, pour tout réel $x$.}
	
	On a sur $\mathbb{R}$, $f'(x) = 3x^2 + 6x + 2$.
	
	\subsubsection*{b. Résoudre dans $\mathbb{R}$ l'équation : $f'(x) = 0$.}
	
	$f'(x) = 0 \iff 3x^2 + 6x + 2 = 0$, \\
	on a $\Delta = 36 - 24 = 12 = (2\sqrt{3})^2 > 0$.
	
	L'équation a donc deux solutions :
	\[
	x_1 = \dfrac{-6 + 2\sqrt{3}}{6} = \dfrac{-3 + \sqrt{3}}{3} \quad \text{et} \quad x_2 = \dfrac{-6 - 2\sqrt{3}}{6} = \dfrac{-3 - \sqrt{3}}{3}
	\]
	
	\subsection*{3. Dresser le tableau de variations de la fonction $f$.}
	
	On sait que $f'(x) > 0$, sauf sur l'intervalle $\left] \dfrac{-3 - \sqrt{3}}{3} ; \dfrac{-3 + \sqrt{3}}{3} \right[$ où $f'(x) < 0$.
	
	Donc $f$ est croissante sauf sur l'intervalle $\left] \dfrac{-3 + \sqrt{3}}{3} ; \dfrac{-3 + \sqrt{3}}{3} \right[$ où elle est décroissante.
	
	\begin{center}
	\begin{tikzpicture}[double distance=2pt]
		\tkzTabInit{$x$/1,$f'(x)$/1,$f(x)$/3}{$-\infty$,$\dfrac{-3 - \sqrt{3}}{3}$,$\dfrac{-3 + \sqrt{3}}{3}$, $+\infty$}
		\tkzTabLine{,+,z,-,z, +}
		\tkzTabVar{-/,+/$1+\dfrac{2\sqrt3}{9}$,-/$1-\dfrac{2\sqrt3}{9}$,+/}
	\end{tikzpicture}
\end{center} 


	
	\subsection*{4. Le point $S(-4 ; -3)$ appartient-il à la tangente à la courbe représentative de $f$ au point d'abscisse $x = -2$ ?}
	
	Équation de la tangente au point $(-2 ; f(-2))$.
	
	On a $f(-2) = (-2)^3 + 3 \times (-2)^2 + 2 \times (-2) + 1 = -8 + 12 - 4 + 1 = 1$.
	
	$f'(-2) = 3 \times (-2)^2 + 6 \times (-2) + 2 = 2$.
	
	La tangente $(T)$ en $S$ a pour équation :
	\begin{align*}
&	M(x ; y) \in (T)\\
	 \iff& y - f(-2) = f'(-2)(x - (-2))\\
	 \iff& y - 1 = 2(x + 2)\\
	 \iff& y = 2x + 4 + 1 \\
	\iff &y = 2x + 5
	\end{align*}
	
	Donc $S(-4 ; -3) \in (T) \iff -3 = 2 \times (-4) + 5 \iff -3 = -8 + 5$ qui est vraie.

\end{document}
